% HanLan (CN++) specification for Babylon Compiler
\documentclass[12pt]{article}
\usepackage[margin=1in]{geometry}
\usepackage{hyperref}
\usepackage{longtable}
\usepackage{amsmath,amssymb}
\usepackage{enumitem}
\usepackage{microtype}

% --- Engine / Font strategy (ULTRA-ROBUST) ---
\usepackage{iftex}
\ifXeTeX\else
  \errmessage{Compile with XeLaTeX in Overleaf (Menu -> Settings -> Compiler -> XeLaTeX).}
\fi
\usepackage{fontspec}
\defaultfontfeatures{Ligatures=TeX}
\setmainfont{FreeSerif}[Renderer=HarfBuzz]
\setmonofont{FreeSerif}[Renderer=HarfBuzz]
\newfontfamily\devanagarifont{FreeSerif}[Script=Devanagari,Renderer=HarfBuzz]
\newcommand{\textdevanagarifont}[1]{{\devanagarifont #1}}

% RFC-style requirement keywords
\newcommand{\MUST}{\textbf{MUST}}
\newcommand{\SHOULD}{\textbf{SHOULD}}
\newcommand{\MAY}{\textbf{MAY}}
\newcommand{\MUSTNOT}{\textbf{MUST NOT}}
\newcommand{\SHOULDNOT}{\textbf{SHOULD NOT}}

\title{HanLan (CN++): Frontend and Registry Specification for Babylon Compiler}
\author{Sebastian Klemm}
\date{\today}

\begin{document}
\maketitle
\tableofcontents

\section{Abstract}
This document specifies \emph{HanLan (CN++)}, a glyph-language designed to operate as an additional
frontend for the \emph{Babylon Compiler} (``Glyph Foundry'' pipeline). The objective is to extend the
existing deterministic compilation pipeline with a new surface language while preserving the
core invariants: Unicode normalization bound in the Run Descriptor (RD), canonical IR lowering,
fixed-point identity math, canonical JSON hashing (RFC 8785), bundle emission, MEF tamper evidence,
and replay verification.

The specification is structured to be implementable directly in the existing Babylon Compiler
workspace as a self-contained addition: new frontend module, CLI selection, test vectors, and
documentation updates.

\section{Repository Baseline}
The Babylon Compiler repository is a Rust workspace comprising crates for canonicalization, IR,
registries, embedding, TIC convergence, crystal construction, MEF chain, gating, macro expansion,
bundle emission, verification, and a CLI pipeline controller. The public interface is the 8-stage
pipeline \texttt{S0..S7} and the \texttt{glyph} CLI.

\subsection{Normative Baseline Invariants}
The following properties are treated as normative baseline invariants for this specification:
\begin{enumerate}[leftmargin=1.5em]
  \item Determinism: identical (normalized source bytes, RD) yields identical IR bytes and manifest digests.
  \item Unicode normalization: \texttt{S0} applies NFC/NFD according to RD; forbidden ranges are enforced.
  \item Single IR: all frontends lower into \texttt{glyph-ir::IrDocument} with canonicalization applied.
  \item Canonical JSON hashing: identity-bearing serialization uses RFC 8785 canonical JSON.
  \item Bundle emission: \texttt{S7} emits a complete bundle (manifest, artifacts, digests, traces).
  \item Verification: \texttt{glyph verify} recomputes digests and validates MEF chain integrity.
\end{enumerate}

\section{HanLan (CN++) Overview}
\subsection{Name, ID, and File Extension}
\begin{itemize}[leftmargin=1.5em]
  \item Language Name: \textbf{HanLan}
  \item Marketing Alias: \textbf{CN++} (``technical Chinese++'' as a naming motif)
  \item Frontend ID (CLI): \texttt{hanlan}
  \item Recommended file extension: \texttt{.hnln}
\end{itemize}

\subsection{Design Objective}
HanLan is a surface language whose \emph{primary differentiator} is its symbol set and operator
lexicon. It is explicitly intended to demonstrate that glyph-based programming surfaces can be
treated as deterministic projections into a canonical IR---not as a replacement of the semantic core.

\section{Lexical Policy}
\subsection{Unicode Policy}
HanLan sources \MUST\ be processed by the existing \texttt{S0 Ingest} stage:
\begin{itemize}[leftmargin=1.5em]
  \item BOM stripping \MUST\ occur before normalization.
  \item Normalization \MUST\ follow RD (\texttt{NFC} default; \texttt{NFD} optional).
  \item Forbidden codepoint ranges \MUST\ be enforced exactly as in the baseline pipeline.
\end{itemize}

\subsection{Identifier Classes}
The lexer \MUST\ accept:
\begin{itemize}[leftmargin=1.5em]
  \item ASCII identifiers: \texttt{[A-Za-z\_][A-Za-z0-9\_]*}
  \item Devanagari identifiers: characters in U+0900..U+097F, including combining marks.
  \item Han identifiers \MAY\ be accepted as identifiers (CJK Unified Ideographs), but \MUST\ be
        normalized and treated deterministically.
\end{itemize}

\subsection{Token Set}
HanLan \MUST\ support the same minimal token and operator set as the Sanskroot MVP, in order to reuse
the existing IR schema and test harness:
\begin{itemize}[leftmargin=1.5em]
  \item Delimiters: \texttt{( ) \{ \} , ;}
  \item Operators: \texttt{+ - * / > < >= <= == != =}
  \item Literals: integer (ASCII digits), string (double quotes), booleans (see below)
\end{itemize}

\section{Grammar (EBNF)}
HanLan grammar is intentionally aligned with the existing Sanskroot AST shape so that the lowering pass
can reuse the same IR node kinds.

\begin{verbatim}
program       ::= { function }

function      ::= FN identifier '(' [ paramlist ] ')' '{' { statement } '}'
paramlist     ::= identifier { ',' identifier }

statement     ::= ifstmt | return | assignment | exprstmt
ifstmt        ::= IF '(' expression ')' '{' { statement } '}'
                 [ ELSE '{' { statement } '}' ]
return        ::= RETURN expression ';'
assignment    ::= LET identifier '=' expression ';'
exprstmt      ::= expression ';'

expression    ::= comparison
comparison    ::= addition [ ( '>' | '<' | '>=' | '<=' | '==' | '!=' ) addition ]
addition      ::= term { ( '+' | '-' ) term }
term          ::= factor { ( '*' | '/' ) factor }
factor        ::= literal | identifier | call | '(' expression ')'
call          ::= PRINT '(' [ arglist ] ')' | identifier '(' [ arglist ] ')'
arglist       ::= expression { ',' expression }
literal       ::= INT | STRING | BOOL
\end{verbatim}

\section{Keyword Vocabulary}
HanLan defines a \emph{CN++} keyword skin. Two keyword tiers are specified:

\subsection{Tier-0 (MVP, Normative)}
Tier-0 keywords are the minimal set required to compile and test:
\begin{longtable}{p{0.22\linewidth}p{0.25\linewidth}p{0.45\linewidth}}
\textbf{Role} & \textbf{Token} & \textbf{Meaning} \\ \hline
Function & \texttt{函} & Function definition keyword (FN) \\
If & \texttt{若} & Conditional keyword (IF) \\
Else & \texttt{否则} & Else branch keyword (ELSE) \\
Return & \texttt{返} & Return keyword (RETURN) \\
Let/Assign & \texttt{令} & Let/assignment keyword (LET) \\
Print & \texttt{示} & Print keyword (PRINT) \\
True/False & \texttt{真}/\texttt{假} & Boolean literals \\
\end{longtable}

\subsection{Tier-1 (Optional Compatibility Layer)}
Tier-1 permits Sanskroot keywords as aliases to allow mixed scripts (useful for debugging and gradual migration):
\begin{itemize}[leftmargin=1.5em]
  \item \textdevanagarifont{कार्य}, \textdevanagarifont{यदि}, \textdevanagarifont{अन्यथा},
        \textdevanagarifont{निवृत्ति}, \textdevanagarifont{मान}, \textdevanagarifont{दर्शय},
        \textdevanagarifont{सत्य}, \textdevanagarifont{असत्य}
\end{itemize}
Tier-1 \MAY\ be disabled by an RD policy flag to enforce pure HanLan sources.

\section{IR Lowering Contract}
\subsection{Canonical IR Target}
HanLan \MUST\ lower into \texttt{glyph-ir::IrDocument}. The document's \texttt{source\_language} field
\MUST\ be \texttt{"hanlan"}.

\subsection{Node/Edge Mapping}
Lowering \MUST\ follow the same mapping discipline as the existing Sanskroot lowering pass:
\begin{itemize}[leftmargin=1.5em]
  \item A root \texttt{Module} node exists with stable content-hash ID.
  \item Each function becomes a \texttt{Function} node with optional \texttt{params}.
  \item Each function body becomes a \texttt{Block} node contained by the function.
  \item Statements map to: \texttt{If}, \texttt{Return}, \texttt{Assignment}, and expression nodes.
  \item Expression nodes map to: \texttt{BinaryOp}, \texttt{Call}, \texttt{Identifier}, \texttt{Literal}.
  \item Ordering edges: \texttt{Contains} edges carry ordinals; \texttt{Next} edges link sequential units.
\end{itemize}
Node IDs \MUST\ be content-hash derived and deterministic; internal counters used for uniqueness \MUST\ be
traversal-order deterministic.

\section{Macro Expansion and CN++ Operator Registry}
\subsection{Macro Registry Integration}
HanLan \MAY\ introduce CN++ ``glyph operators'' as macro expansions executed in \texttt{S2 Expand}.
All macro expansion \MUST\ be deterministic given (normalized source, RD, macro registry digest).

\subsection{Non-Normative Example Macros}
The following expansions illustrate the intended style:
\begin{itemize}[leftmargin=1.5em]
  \item \texttt{火(x)} $\Rightarrow$ \texttt{示(x)} (print)
  \item \texttt{门(c,a,b)} $\Rightarrow$ \texttt{若(c)\{a;\}否则\{b;\}}
\end{itemize}

\section{Integration into Babylon Compiler Workspace}
\subsection{Required Code Additions (File-Level)}
The implementation \MUST\ add a new frontend module in \texttt{glyph-frontends}:
\begin{itemize}[leftmargin=1.5em]
  \item \texttt{crates/glyph-frontends/src/hanlan/mod.rs}
  \item \texttt{crates/glyph-frontends/src/hanlan/lexer.rs}
  \item \texttt{crates/glyph-frontends/src/hanlan/parser.rs}
  \item \texttt{crates/glyph-frontends/src/hanlan/lower.rs}
  \item Update \texttt{crates/glyph-frontends/src/lib.rs} to export \texttt{pub mod hanlan;}
\end{itemize}

\subsection{CLI Frontend Selection}
The CLI stage \texttt{S1 Canon} \MUST\ accept \texttt{--frontend hanlan} and route to the HanLan frontend.

\subsection{Documentation Updates}
\begin{itemize}[leftmargin=1.5em]
  \item \texttt{README.md} \SHOULD\ document HanLan, its keywords, and \texttt{.hnln} extension.
  \item A new file \texttt{examples/hello.hnln} \SHOULD\ be added as a canonical example.
\end{itemize}

\section{Test Vectors (Normative)}
Three canonical test vectors \MUST\ be added under \texttt{glyph-frontends}.

\subsection{TV1: Minimal Print}
\begin{verbatim}
函 主() {
  示("你好");
}
\end{verbatim}

\subsection{TV2: Conditional}
\begin{verbatim}
函 检(x) {
  若 (x > 0) {
    示("正");
  } 否则 {
    示("非正");
  }
}
\end{verbatim}

\subsection{TV3: Multi-Function}
\begin{verbatim}
函 和(a, b) { 返 a + b; }
函 主() {
  令 r = 和(3, 4);
  示(r);
}
\end{verbatim}

\subsection{Determinism Assertion}
A unit test \MUST\ assert that lowering the same HanLan source twice yields identical document digests.

\section{Acceptance Criteria (Definition of Done)}
An implementation satisfies this specification iff:
\begin{enumerate}[leftmargin=1.5em]
  \item \texttt{glyph run --frontend hanlan} compiles TV1--TV3 and emits a bundle.
  \item The bundle passes \texttt{glyph verify} without errors.
  \item \texttt{glyph dump-ir --frontend hanlan} emits canonical JSON for the IR.
  \item Determinism: repeated runs yield identical IR digest and manifest digest under identical RD and source.
  \item Documentation includes HanLan keyword table and a \texttt{.hnln} example.
\end{enumerate}

\end{document}
